\addcontentsline{toc}{subsubsection}{Solution 43. The Moving Zero}

\begin{dlsummary}[title={\bfseries Solution 43. The Moving Zero},breakable]

\begin{proof}
The only solution is
\[
\boxed{f(x)=x^2\qquad(x\in\mathbb R).}
\]

\paragraph{1. Find a zero—but do not assume it is \(0\).}

Set
\[
a=f(0)\ge0.
\]
By (2) at \(x=0\),
\[
0\le f(a)=f(f(0))\le0,
\]
so \(f(a)=0\).

\paragraph{2. Build a quadratic lower bound around that zero.}

In (1), substitute
\[
x=\frac{t+a}{2},\qquad y=\frac{t-a}{2}.
\]
Since \(f(a)=0\),
\[
f(t)\ge 2f\!\left(\frac{t+a}{2}\right)+\frac{(t-a)^2}{2}.
\]
Repeating this midpoint substitution \(n\) times gives
\[
f(t)\ge
2^n f\!\left(a+\frac{t-a}{2^n}\right)
+(1-2^{-n})(t-a)^2.
\]
Nonnegativity lets us discard the first term. Letting \(n\to\infty\), we obtain
\[
\boxed{f(t)\ge(t-a)^2\qquad(t\in\mathbb R).} \tag{3}
\]
No limit of function values is being taken, so continuity is unnecessary.

\paragraph{3. Turn the lower bound into an upper bound.}

Apply (3) at \(t=f(x)\), then use (2):
\[
(f(x)-a)^2\le f(f(x))\le x^4.
\]
Therefore
\[
|f(x)-a|\le x^2,
\qquad\text{in particular}\qquad
f(x)\le a+x^2. \tag{4}
\]

On the other hand, (1) at \(x=0\) gives
\[
f(y)+f(-y)\ge2a+2y^2.
\]
Each summand is at most \(a+y^2\) by (4). Both must therefore attain their upper bounds:
\[
f(y)=a+y^2\qquad(y\in\mathbb R).
\]
Since \(f(a)=0\),
\[
a+a^2=0.
\]
Together with \(a\ge0\), this forces \(a=0\). Hence \(f(x)=x^2\).

Finally, \(f(x)=x^2\) satisfies both original conditions with equality.
\end{proof}

\paragraph{Takeaways.}
\begin{itemize}
\item **Locate a zero before normalising:** The composition condition supplies a zero at \(f(0)\), not immediately at \(0\).
\item **Iterate a midpoint inequality:** Repeated bisection produces a sharp bound without establishing continuity.
\item **Feed bounds into compositions:** A lower bound for \(f(t)\) becomes an upper bound for \(f(x)\) through \(f(f(x))\).
\item **Exploit equality of sums:** If two terms each have the same upper bound and their sum reaches the total bound, both are determined.
\end{itemize}

\end{dlsummary}
